 \lecture
{07}
{Friday, 28 August 2026}
{Introduction to Probability}
{Dr. Nishant Chandgotia}

%%%%%%%%%%%%%%%%%%%%%%%%%%%%%%%%%%%%%%%%%%%%%%%%%%%%%%%%%%%%
\section{Examples of Random Variables}
%%%%%%%%%%%%%%%%%%%%%%%%%%%%%%%%%%%%%%%%%%%%%%%%%%%%%%%%%%%%

\begin{example}{Uniform Random Variable}{uniform-random-variable}

Let $X$ be uniformly distributed on $(0,1)$. Then its C.D.F. is
given by
\[
F_X(x)
=
\begin{cases}
0, & x\leq 0,\\[4pt]
x, & 0<x<1,\\[4pt]
1, & x\geq 1.
\end{cases}
\]

\end{example}


\begin{example}{Exponential Random Variable}{exponential-random-variable}

Let
\[
X\sim\operatorname{Exp}(\lambda),
\qquad \lambda>0.
\]
Then $X$ has density
\[
f_X(y)=\lambda e^{-\lambda y},
\qquad y\ge 0.
\]

\end{example}


\begin{example}{Poisson Random Variable}{poisson-random-variable}

Let
\[
X\sim\operatorname{Poisson}(\lambda),
\qquad \lambda>0.
\]
Then the C.D.F. of $X$ is
\[
F_X(x)
=
\begin{cases}
0, & x<0,\\[6pt]
\displaystyle
\sum_{k=0}^{\lfloor x\rfloor}
e^{-\lambda}\frac{\lambda^k}{k!},
& x\geq 0.
\end{cases}
\]

\end{example}


\begin{example}{Standard Normal Random Variable}{standard-normal-random-variable}

Let
\[
X\sim N(0,1).
\]
Then the C.D.F. of $X$ is
\[
F_X(x)
=
\int_{-\infty}^{x}
\frac{1}{\sqrt{2\pi}}
e^{-y^2/2}\,dy.
\]

\end{example}


%%%%%%%%%%%%%%%%%%%%%%%%%%%%%%%%%%%%%%%%%%%%%%%%%%%%%%%%%%%%
\section{A Gaussian Tail Estimate}
%%%%%%%%%%%%%%%%%%%%%%%%%%%%%%%%%%%%%%%%%%%%%%%%%%%%%%%%%%%%

\begin{theorem}{Gaussian Tail Bound}{gaussian-tail-bound}

For $x>0$,
\[
\left(
\frac{1}{x}-\frac{1}{x^3}
\right)e^{-x^2/2}
\leq
\int_x^\infty e^{-y^2/2}\,dy
\leq
\frac{1}{x}e^{-x^2/2}.
\]

\end{theorem}


\begin{proof}

We first prove the upper bound.

Let
\[
y=x+z.
\]
Then
\[
\begin{aligned}
\int_x^\infty e^{-y^2/2}\,dy
&=
e^{-x^2/2}
\int_0^\infty
e^{-xz-z^2/2}\,dz.
\end{aligned}
\]

Since
\[
e^{-xz-z^2/2}
\leq e^{-xz},
\]
we obtain
\[
\begin{aligned}
\int_x^\infty e^{-y^2/2}\,dy
&\leq
e^{-x^2/2}
\int_0^\infty e^{-xz}\,dz\\
&=
\frac{1}{x}e^{-x^2/2}.
\end{aligned}
\]

Thus,
\[
\boxed{
\int_x^\infty e^{-y^2/2}\,dy
\leq
\frac{1}{x}e^{-x^2/2}
}.
\]

\end{proof}
Todo: We have to compute the lower bound also.

%%%%%%%%%%%%%%%%%%%%%%%%%%%%%%%%%%%%%%%%%%%%%%%%%%%%%%%%%%%%
\section{Expectation and Variance}
%%%%%%%%%%%%%%%%%%%%%%%%%%%%%%%%%%%%%%%%%%%%%%%%%%%%%%%%%%%%

\begin{definition}{Expectation}{expectation}

Suppose $X$ is a random variable such that
\[
X\in L^1(\mathbb{P}).
\]
Then the expectation of $X$ is defined by
\[
\mathbb{E}(X)
=
\int_\Omega X\,d\mathbb{P}.
\]

\end{definition}


\begin{definition}{Variance}{variance}

Suppose
\[
X\in L^2(\mathbb{P}).
\]
Then the variance of $X$ is defined by
\[
\operatorname{Var}(X)
=
\mathbb{E}
\left[
\left(X-\mathbb{E}(X)\right)^2
\right].
\]

\end{definition}


%%%%%%%%%%%%%%%%%%%%%%%%%%%%%%%%%%%%%%%%%%%%%%%%%%%%%%%%%%%%
\section{Central Limit Theorem}
%%%%%%%%%%%%%%%%%%%%%%%%%%%%%%%%%%%%%%%%%%%%%%%%%%%%%%%%%%%%

\begin{theorem}{Central Limit Theorem}{central-limit-theorem}

Let
\(
X_1,X_2,\ldots
\)
be i.i.d.\ random variables such that they  all have same distribution say X,~ ~ ~ 
\(
\mathbb{E}(X_1^2)<\infty.
\)
Let
\[
\mu=\mathbb{E}(X_1),
\qquad
\sigma^2=\operatorname{Var}(X_1).
\]

Then
\[
\frac{
\displaystyle
\sum_{i=1}^n X_i-n\mu
}{
\sqrt{n}\sigma
}
\xrightarrow{d}
N(0,1).
\]

Equivalently, for $a<b$,
\[
\mathbb{P}
\left(
a
\leq
\frac{
\displaystyle\sum_{i=1}^n X_i-n\mu
}{
\sqrt{n}\sigma
}
\leq b
\right)
\longrightarrow
\frac{1}{\sqrt{2\pi}}
\int_a^b e^{-x^2/2}\,dx.
\]

\end{theorem}

\begin{remark}
The result above is the usual form of the Central Limit Theorem.
\end{remark}


%%%%%%%%%%%%%%%%%%%%%%%%%%%%%%%%%%%%%%%%%%%%%%%%%%%%%%%%%%%%
\section{Examples of Expectation and Variance}
%%%%%%%%%%%%%%%%%%%%%%%%%%%%%%%%%%%%%%%%%%%%%%%%%%%%%%%%%%%%

\begin{example}{Expectation and Variance of Common Distributions}
{expectation-variance-common-distributions}

\begin{enumerate}

    \item If
    \(
    X\sim\operatorname{Exp}(\lambda),
    \)
    then
    \[
    \mathbb{E}(X)=\frac{1}{\lambda},
    \qquad
    \operatorname{Var}(X)=\frac{1}{\lambda^2}.
    \]

    \item If
    \(
    X\sim\operatorname{Poisson}(\lambda),
    \)
    then
    \[
    \mathbb{E}(X)=\lambda,
    \qquad
    \operatorname{Var}(X)=\lambda.
    \]

    \item If
    \(
    X\sim\operatorname{Geom}(p),
    \)
    then
    \[
    \mathbb{E}(X)=\frac{1}{p},
    \qquad
    \operatorname{Var}(X)
    =
    \frac{1-p}{p^2}.
    \]

    \item If
    \(
    X\sim N(\mu,\sigma^2),
    \)
    then
    \[
    \mathbb{E}(X)=\mu,
    \qquad
    \operatorname{Var}(X)=\sigma^2.
    \]

\end{enumerate}

\end{example}


\begin{question}{Prove this}{prove-t}
Do the calculations to prove
\(
\mathbb{E}(X)=\mu
\qquad\text{and}\qquad
\operatorname{Var}(X)=\sigma^2
\)
for
\(
X\sim N(\mu,\sigma^2).
\)
\end{question}


\begin{remark}

If
\(
Y=aX+b,
\)
then
\(
\mathbb{E}(Y)
=
a\mathbb{E}(X)+b
\)
and
\(
\operatorname{Var}(Y)
=
a^2\operatorname{Var}(X).
\)

\end{remark}


%%%%%%%%%%%%%%%%%%%%%%%%%%%%%%%%%%%%%%%%%%%%%%%%%%%%%%%%%%%%
\section{A Result on Signs of Vectors}
%%%%%%%%%%%%%%%%%%%%%%%%%%%%%%%%%%%%%%%%%%%%%%%%%%%%%%%%%%%%

\begin{theorem}{Sign Choice for Vectors}{sign-choice-vectors}

Let
\(
v_1,\ldots,v_n\in\mathbb{R}^n
\)
satisfy
\(
\|v_i\|_2=1
\qquad
\text{for all }i.
\)

Then:

\begin{enumerate}

    \item There exist
    \(
    \varepsilon_1,\ldots,\varepsilon_n
    \in\{-1,1\}
    \)
    such that
    \(
    \left\|
    \sum_{i=1}^n
    \varepsilon_i v_i
    \right\|_2
    \leq
    \sqrt{n}.
    \)

    \item There exist
    \(
    \varepsilon_1,\ldots,\varepsilon_n
    \in\{-1,1\}
    \)
    such that
    \(
    \left\|
    \sum_{i=1}^n
    \varepsilon_i v_i
    \right\|_2
    \geq
    \sqrt{n}.
    \)

\end{enumerate}

\end{theorem}


\begin{proof}

Let
\(
\varepsilon_1,\ldots,\varepsilon_n
\)
be independent random variables, each uniformly distributed on
$\{-1,1\}$.

Then
\[
\begin{aligned}
\mathbb{E}
\left[
\left\|
\sum_{i=1}^n\varepsilon_i v_i
\right\|_2^2
\right]
&=
\mathbb{E}
\left[
\sum_{i,j}
\varepsilon_i\varepsilon_j
\langle v_i,v_j\rangle
\right]\\
&=
\sum_{i,j}
\langle v_i,v_j\rangle
\mathbb{E}(\varepsilon_i\varepsilon_j).
\end{aligned}
\]

For $i\neq j$,
\(
\mathbb{E}(\varepsilon_i\varepsilon_j)
=
\mathbb{E}(\varepsilon_i)
\mathbb{E}(\varepsilon_j)
=
0,
\)
while
\(
\mathbb{E}(\varepsilon_i^2)=1.
\)

Therefore,
\[
\begin{aligned}
\mathbb{E}
\left[
\left\|
\sum_{i=1}^n\varepsilon_i v_i
\right\|_2^2
\right]
&=
\sum_{i=1}^n
\|v_i\|_2^2\\
&=
n.
\end{aligned}
\]

Hence there must exist a choice of signs for which
\(
\left\|
\sum_{i=1}^n\varepsilon_i v_i
\right\|_2^2
\geq n,
\)
and there must also exist a choice of signs for which
\(
\left\|
\sum_{i=1}^n\varepsilon_i v_i
\right\|_2^2
\leq n.
\)

Taking square roots gives the result.

\end{proof}


%%%%%%%%%%%%%%%%%%%%%%%%%%%%%%%%%%%%%%%%%%%%%%%%%%%%%%%%%%%%
\section{Exercises on Probability Distributions}
%%%%%%%%%%%%%%%%%%%%%%%%%%%%%%%%%%%%%%%%%%%%%%%%%%%%%%%%%%%%

\begin{question}{}{question-1}
Suppose $X$ is a discrete random variable such that
\[
\mathbb{P}(X=k)
=
(1-p)^{k-1}p,
\qquad
k\geq 1,
\qquad
k\in\mathbb{N}.
\]

Find
\[
\mathbb{P}(a<X<b).
\]
\end{question}


\begin{question}{}{question-2}
Suppose $X$ is a random variable whose probability density
function is
\[
f_X(x)
=
\lambda e^{-\lambda x},
\qquad
x\geq0,
\qquad
\lambda>0.
\]

Find
\[
\mathbb{P}(a<X<b).
\]
\end{question}


%%%%%%%%%%%%%%%%%%%%%%%%%%%%%%%%%%%%%%%%%%%%%%%%%%%%%%%%%%%%
\section{Graph Theory}
%%%%%%%%%%%%%%%%%%%%%%%%%%%%%%%%%%%%%%%%%%%%%%%%%%%%%%%%%%%%

\begin{definition}{Graph}{graph}

A graph is a pair
\(
G=(V,E),
\)
where

\begin{itemize}
    \item $V$ is the set of vertices;
    \item $E$ is the set of edges.
\end{itemize}

In this lecture, we consider undirected graphs.

\end{definition}


\begin{example}{Basic Graphs}{basic-graphs}

Examples of graphs include:

\begin{enumerate}

    \item The complete simple graph $K_n$.

    \item The cycle graph $C_n$, where
    \(
    V=\mathbb{Z}_n
    \)
    and
    \(
    E
    =
    \left\{
    (i,i+1):i\in\mathbb{Z}_n
    \right\}.
    \)

\end{enumerate}

\end{example}


%%%%%%%%%%%%%%%%%%%%%%%%%%%%%%%%%%%%%%%%%%%%%%%%%%%%%%%%%%%%
\section{Bipartite Graphs}
%%%%%%%%%%%%%%%%%%%%%%%%%%%%%%%%%%%%%%%%%%%%%%%%%%%%%%%%%%%%

\begin{definition}{Bipartite Graph}{bipartite-graph}

A graph $G=(V,E)$ is called bipartite if there exists a partition
\(
V=V_1\sqcup V_2
\)
such that whenever
\(
(u,v)\in E,
\)
we have
\(
u\in V_1,\qquad v\in V_2,
\)
or equivalently,
\(
u\in V_2,\qquad v\in V_1.
\)

\end{definition}


\begin{example}{Odd Cycle is Not Bipartite}{odd-cycle-not-bipartite}

The cycle graph $C_n$ is bipartite if and only if $n$ is even.

For an odd cycle, one cannot partition the vertices into two sets such
that every edge joins vertices belonging to different sets.

\end{example}


\begin{example}{Path Graph}{path-graph}

The path graph $P_n$ is always bipartite.

\end{example}


%%%%%%%%%%%%%%%%%%%%%%%%%%%%%%%%%%%%%%%%%%%%%%%%%%%%%%%%%%%%
\section{Subgraphs and Walks}
%%%%%%%%%%%%%%%%%%%%%%%%%%%%%%%%%%%%%%%%%%%%%%%%%%%%%%%%%%%%

\begin{definition}{Subgraph}{subgraph}

A graph
\(
G'=(V',E')
\)
is a subgraph of
\(
G=(V,E)
\)
if
\(
V'\subseteq V
\qquad\text{and}\qquad
E'\subseteq E.
\)

\end{definition}


\begin{definition}{Walk}{walk}

A walk from a vertex to another vertex in
\(
G=(V,E)
\)
is a sequence of vertices
\(
v_0,v_1,\ldots,v_n
\)
such that
\(
(v_i,v_{i+1})\in E
\)
for every
\(
i=0,\ldots,n-1.
\)

\end{definition}


\begin{definition}{Tree}{tree}

A graph is called a tree if it is connected and does not contain a
cycle as a subgraph.

\end{definition}


\begin{definition}{Connected Graph}{connected-graph}

A graph is called connected if, for all distinct vertices
$u,v\in V$, there exists a walk from $u$ to $v$.

\end{definition}


%%%%%%%%%%%%%%%%%%%%%%%%%%%%%%%%%%%%%%%%%%%%%%%%%%%%%%%%%%%%
\section{Trees and Bipartite Graphs}
%%%%%%%%%%%%%%%%%%%%%%%%%%%%%%%%%%%%%%%%%%%%%%%%%%%%%%%%%%%%

\begin{theorem}{Every Tree is Bipartite}{every-tree-is-bipartite}

Every tree is a bipartite graph.

\end{theorem}


\begin{proof}

Let
\(
T=(V,E)
\)
be a tree and choose a vertex
\(
v\in V.
\)

Define
\[
V_0
=
\left\{
u\in V:
d(u,v)\text{ is even}
\right\}
\]
and
\[
V_1
=
\left\{
u\in V:
d(u,v)\text{ is odd}
\right\},
\]
where $d(u,v)$ denotes the length of the shortest walk from $u$ to
$v$.

Clearly,
\[
V=V_0\sqcup V_1.
\]

Suppose $(u,w)\in E$. If $d(u,v)$ and $d(w,v)$ had the same parity,
then combining shortest paths from $v$ to $u$ and $v$ to $w$ with the
edge $(u,w)$ would produce cycle.

Since $T$ is a tree, this is impossible.

Therefore every edge joins a vertex in $V_0$ to a vertex in $V_1$.
Hence $T$ is bipartite.

\end{proof}


%%%%%%%%%%%%%%%%%%%%%%%%%%%%%%%%%%%%%%%%%%%%%%%%%%%%%%%%%%%%
\section{Characterization of Bipartite Graphs}
%%%%%%%%%%%%%%%%%%%%%%%%%%%%%%%%%%%%%%%%%%%%%%%%%%%%%%%%%%%%

\begin{theorem}{Characterization of Bipartite Graphs}
{bipartite-odd-cycle-characterization}

A graph
\(
G=(V,E)
\)
is bipartite if and only if it contains no cycle of odd length.

\end{theorem}


\begin{remark}

An odd cycle is not bipartite.

Moreover, every subgraph of a bipartite graph is also bipartite.

\end{remark}


\begin{proof}

We have already seen that every tree is bipartite.

Now suppose that $G$ is bipartite. Let
\(
V=V_0\sqcup V_1
\)
be a bipartition of $V$.

Consider a cycle
\[
v_0,v_1,\ldots,v_{n-1},v_0.
\]

Since every edge joins the two different parts of the bipartition,
the vertices of the cycle must alternate between $V_0$ and $V_1$.

Therefore the length of the cycle must be even.

Hence every cycle in a bipartite graph is even.

\end{proof}


%%%%%%%%%%%%%%%%%%%%%%%%%%%%%%%%%%%%%%%%%%%%%%%%%%%%%%%%%%%%
\section{Bipartite Subgraph}
%%%%%%%%%%%%%%%%%%%%%%%%%%%%%%%%%%%%%%%%%%%%%%%%%%%%%%%%%%%%

\begin{theorem}{Bipartite Subgraph}{large-bipartite-subgraph}

Let
\(
G=(V,E)
\)
be a finite graph without self-loops. Then there exists a bipartite
subgraph
\(
G'=(V,E')
\)
such that
\(
|E'|\geq\frac{|E|}{2}.
\)

\end{theorem}


\begin{proof}

Let $V$ be the vertex set and consider a random partition
\(
V=V_0\sqcup V_1
\)
obtained by independently assigning each vertex to $V_0$ or $V_1$ with
probability $1/2$.

An edge $(u,v)\in E$ is said to be \emph{crossing} if $u$ and $v$
belong to different parts of the partition.

Define the random variable
\[
X_{(u,v)}
=
\begin{cases}
1, & \text{if $(u,v)$ is crossing},\\
0, & \text{otherwise}.
\end{cases}
\]

Let
\(
X
=
\sum_{(u,v)\in E}
X_{(u,v})
\)
be the total number of crossing edges.

For a fixed edge $(u,v)$,
\[
\begin{aligned}
\mathbb{E}(X_{(u,v)})
&=
\mathbb{P}(X_{(u,v)}=1)\\
&=
\mathbb{P}(u\in V_0,v\in V_1)
+
\mathbb{P}(u\in V_1,v\in V_0)\\
&=
\frac14+\frac14\\
&=
\frac12.
\end{aligned}
\]

Therefore,
\[
\begin{aligned}
\mathbb{E}(X)
&=
\sum_{(u,v)\in E}
\mathbb{E}(X_{(u,v)})\\
&=
\frac{|E|}{2}.
\end{aligned}
\]

Hence there must exist at least one partition for which
\[
X\geq\frac{|E|}{2}.
\]

Take $E'$ to be the set of crossing edges for such a partition.
Then
\(
G'=(V,E')
\)
is bipartite and satisfies
\(
|E'|\geq\frac{|E|}{2}.
\)

\end{proof}


%%%%%%%%%%%%%%%%%%%%%%%%%%%%%%%%%%%%%%%%%%%%%%%%%%%%%%%%%%%%
%% End of Lecture
%%%%%%%%%%%%%%%%%%%%%%%%%%%%%%%%%%%%%%%%%%%%%%%%%%%%%%%%%%%%

\lectureend
